\documentclass[10pt,twocolumn,letterpaper]{article}

%% ── Packages ──────────────────────────────────────────────────────────────
\usepackage[utf8]{inputenc}
\usepackage[T1]{fontenc}
\usepackage{lmodern}
\usepackage[margin=0.75in]{geometry}
\usepackage{amsmath,amssymb}
\usepackage{booktabs}
\usepackage{graphicx}
\usepackage{xcolor}
\usepackage{hyperref}
\usepackage{enumitem}
\usepackage{titlesec}
\usepackage{caption}
\usepackage{subcaption}
\usepackage{fancyhdr}
\usepackage{listings}
\usepackage{algorithm}
\usepackage{algpseudocode}
\usepackage{tabularx}
\usepackage{float}

%% ── Colors ────────────────────────────────────────────────────────────────
\definecolor{cascadeblue}{HTML}{0088FF}
\definecolor{cascadedark}{HTML}{001824}
\definecolor{codebg}{HTML}{F5F5F5}
\hypersetup{
  colorlinks=true,
  linkcolor=cascadeblue,
  urlcolor=cascadeblue,
  citecolor=cascadeblue
}

%% ── Listings ──────────────────────────────────────────────────────────────
\lstset{
  backgroundcolor=\color{codebg},
  basicstyle=\ttfamily\scriptsize,
  breaklines=true,
  frame=single,
  rulecolor=\color{gray!30},
  columns=fullflexible,
  keepspaces=true,
  xleftmargin=4pt,
  xrightmargin=4pt,
  aboveskip=6pt,
  belowskip=6pt
}

%% ── Section formatting ───────────────────────────────────────────────────
\titleformat{\section}{\large\bfseries}{\thesection.}{0.5em}{}
\titleformat{\subsection}{\normalsize\bfseries}{\thesubsection}{0.5em}{}
\titlespacing*{\section}{0pt}{10pt}{4pt}
\titlespacing*{\subsection}{0pt}{8pt}{3pt}
\setlength{\columnsep}{0.25in}

%% ── Header/Footer ────────────────────────────────────────────────────────
\pagestyle{fancy}
\fancyhf{}
\fancyhead[L]{\small\textsc{CASCADE}}
\fancyhead[R]{\small National Space Hackathon 2026}
\fancyfoot[C]{\thepage}
\renewcommand{\headrulewidth}{0.4pt}

%% ── Title ─────────────────────────────────────────────────────────────────
\title{
  \vspace{-1.5em}
  {\LARGE\bfseries CASCADE: Autonomous Orbital Collision\\[2pt]
  Avoidance for Congested LEO}\\[6pt]
  {\normalsize National Space Hackathon 2026 --- IIT Delhi}
}
\author{
  Team Space\_Invaders\\[2pt]
  \small \url{https://github.com/Aqwerty321/PROJECT-B.O.N.K}
}
\date{March 2026}

\begin{document}
\maketitle
\thispagestyle{fancy}

%% ═══════════════════════════════════════════════════════════════════════════
\begin{abstract}
\noindent
Low Earth Orbit faces an escalating Kessler-syndrome risk as mega-constellation
deployments collide with a legacy debris population exceeding 27{,}000 tracked
objects.  We present \textbf{CASCADE}, a deterministic, single-container
Autonomous Constellation Manager that ingests high-frequency ECI telemetry,
screens for conjunctions through a multi-tier filtering cascade, and
autonomously schedules collision-avoidance and recovery burns under fuel,
cooldown, and line-of-sight constraints.  The engine processes 50~satellites
against 10{,}000~debris fragments in \textbf{13.3\,ms} mean tick latency with
\textbf{zero false negatives} across all test families.  A React/Canvas/WebGL
frontend delivers a mission-console experience at 60\,FPS.  CASCADE ships as a
single Docker image (\texttt{ubuntu:22.04}, port~8000) with 8/8~CTest safety
gates enforced in CI, meeting every deployment and API contract specified in
the problem statement.
\end{abstract}

%% ═══════════════════════════════════════════════════════════════════════════
\section{Introduction and Design Goals}

The National Space Hackathon 2026 challenges teams to build an Autonomous
Constellation Manager (ACM) that can (i)~ingest high-frequency ECI state
vectors, (ii)~predict conjunctions up to 24\,hours ahead without $O(N^{2})$
bottlenecks, (iii)~autonomously compute and schedule evasion and recovery
maneuvers, and (iv)~visualize the operational state in a real-time mission
dashboard.

CASCADE addresses these objectives under a strict \emph{safety-first} design
philosophy:

\begin{itemize}[nosep,leftmargin=*]
  \item \textbf{Zero false negatives}: every ambiguous pair is escalated to a
    more expensive solver, never silently dropped.
  \item \textbf{Deterministic runtime}: identical inputs always produce the same
    maneuver sequence, enabling reproducible CI gates.
  \item \textbf{Conservative classification}: the engine prefers false positives
    (unnecessary screening work) over false negatives (missed collisions).
  \item \textbf{Single-container deployment}: one \texttt{docker build} yields a
    self-contained binary and frontend, exposed on port~8000.
\end{itemize}

The remainder of this report details the system architecture
(\S\ref{sec:arch}), key design decisions and trade-offs
(\S\ref{sec:decisions}), numerical methods (\S\ref{sec:methods}), frontend
console (\S\ref{sec:frontend}), validation evidence (\S\ref{sec:validation}),
and deployment pipeline (\S\ref{sec:deploy}).

%% ═══════════════════════════════════════════════════════════════════════════
\section{System Architecture}\label{sec:arch}

Figure~\ref{fig:pipeline} illustrates the end-to-end data path.  The system
is a single-process C++20 engine linked against two FetchContent dependencies
--- \texttt{cpp-httplib}~v0.15.3 (HTTP) and \texttt{simdjson}~v3.9.4 (JSON
parsing) --- with Boost headers for version verification.

\begin{figure}[h]
\centering
\fbox{\parbox{0.95\linewidth}{\small\ttfamily
\begin{tabbing}
POST /api/telemetry \\
\quad $\downarrow$ \\
{[State Store]} -- lock-free SoA upsert \\
\quad $\downarrow$ \\
{[Propagation]} -- adaptive RK4+J2 \\
\quad $\downarrow$ \\
{[Broad Phase]} -- $O(N \log N)$ SMA+incl bins \\
\quad $\downarrow$ \\
{[Narrow Phase]} -- MOID + TCA + RK4 refine \\
\quad $\downarrow$ \\
{[CDM Engine]} -- 24\,h conjunction horizon \\
\quad $\downarrow$ \\
{[COLA Planner]} -- RTN evasion + Tsiolkovsky \\
\quad $\downarrow$ \\
{[Recovery]} -- CW/ZEM slot-targeting burn \\
\quad $\downarrow$ \\
REST API (port 8000)
\end{tabbing}
}}
\caption{End-to-end simulation tick pipeline.}
\label{fig:pipeline}
\end{figure}

\subsection{State Representation}

All objects are stored in a Structure-of-Arrays (SoA) layout indexed by
string identifier.  Each record holds:

\begin{itemize}[nosep,leftmargin=*]
  \item 6-DOF ECI state vector $(\vec{r}, \vec{v})$ in km and km/s
  \item Precomputed Keplerian elements $(a, e, i, \Omega, \omega, M)$
  \item Radial bounds $(r_{\min}, r_{\max})$ for shell overlap
  \item Fuel mass, cooldown timer, slot reference, burn history
\end{itemize}

\subsection{API Surface}

The engine exposes five primary endpoints on port~8000, bound to
\texttt{0.0.0.0}:

\begin{table}[h]
\centering\small
\begin{tabularx}{\linewidth}{@{}lX@{}}
\toprule
\textbf{Endpoint} & \textbf{Purpose} \\
\midrule
\texttt{POST /api/telemetry}       & Bulk ECI state ingestion \\
\texttt{POST /api/maneuver/schedule} & Evasion/recovery burn submission \\
\texttt{POST /api/simulate/step}    & Advance simulation by $\Delta t$\,s \\
\texttt{GET /api/visualization/snapshot} & Geodetic snapshot for frontend \\
\texttt{GET /api/status}            & Engine health \& diagnostics \\
\bottomrule
\end{tabularx}
\caption{Primary REST API endpoints.}
\label{tab:api}
\end{table}

\noindent
Additional debug endpoints (\texttt{/api/debug/\{burns, conjunctions,
propagation, conflicts\}}) expose internal ring buffers for observability.

%% ═══════════════════════════════════════════════════════════════════════════
\section{Design Decisions and Trade-offs}\label{sec:decisions}

Several non-obvious architectural choices emerged from iterative development
over 60+ commits.  We document the key trade-offs to explain why the system
takes its current form.

\subsection{Julia Bridge $\to$ Pure C++ Runtime}

The architecture document originally specified a Julia/\texttt{jluna} hybrid
for heavy numerics (\texttt{DifferentialEquations.jl}, analytical MOID).  We
implemented a complete NSGA-II parameter tuner via \texttt{jluna}
(\texttt{tuner/nsga2\_tuner.jl}) and verified end-to-end optimization over
69~parameters.  However, the Julia~1.10 runtime added \textbf{$\sim$200\,MB}
to the Docker image and introduced JIT latency on container boot.  Since the
grading VM runs a cold \texttt{docker build}, we removed the Julia runtime
from the production Dockerfile and re-implemented all numerical kernels in
pure C++20.  The NSGA-II tuner remains available as an offline tool
(\texttt{PROJECTBONK\_ENABLE\_JULIA\_RUNTIME} CMake guard) but the
submission binary has zero Julia dependency.

\subsection{Shadow-First Filter Rollout}

New narrow-phase filters (plane/phase gate, MOID proxy) are deployed in
\emph{shadow mode} by default: the gate evaluates and logs rejection
statistics but does not actually discard pairs.  This allows us to
calibrate false-negative risk against the full RK4 reference before
activating hard filtering.  The pattern was motivated by an early
observation that the D-criterion gate's threshold was sensitive to orbital
element quality --- activating it prematurely on poorly-determined debris
could violate the zero-false-negative invariant.  Shadow-first rollout
ensures every filter accumulates calibration evidence (exposed via
\texttt{/api/status?details=1}) before moving to production.

\subsection{Heuristic $\to$ CW/ZEM Recovery Solver}

The initial recovery planner used proportional RTN gains
(\texttt{HEURISTIC} mode) --- simple but producing suboptimal $\Delta V$
vectors, especially on eccentric orbits where velocity-direction diverges
from true tangential.  An audit identified three correctness issues:
(i)~mean motion computed from instantaneous radius instead of semi-major
axis, (ii)~$\hat{t}$ approximated as $\hat{v}$ instead of
$\hat{n} \times \hat{r}$, and (iii)~simplified linear CW terms instead of
true $\sin(n\tau)/\cos(n\tau)$ state-transition matrix.  All three were
fixed, and \texttt{CW\_ZEM\_EQUIVALENT} became the default solver, with
the heuristic retained as a near-singularity fallback.

\subsection{3D Globe $\to$ 2D Operations Dashboard}

The frontend initially featured a full-viewport Three.js globe with
Fresnel atmosphere, glassmorphic HUD panels, and a ``TRON $\times$ Star
Wars'' cockpit aesthetic.  While visually impressive, judge feedback
priority analysis revealed that the PS requires a \emph{2D operational
dashboard} with ground tracks, maneuver timelines, and resource heatmaps.
We pivoted: the 3D globe was demoted to a supporting context element,
and a Canvas-based Mercator ground track with 90-minute trail/forecast,
live terminator, and a scrollable Gantt maneuver timeline became the
primary surfaces.  This decision improved judge-rubric alignment
(Section~6 compliance) at the cost of visual spectacle.

\subsection{MOID: Sampled Proxy vs.\ Analytical Solver}

The problem statement rewards algorithmic sophistication, so we initially
targeted a Gronchi~(2005) analytical quartic MOID solver.  We implemented
it (\texttt{orbit\_math.cpp}) and observed correct results, but the
analytical path's sensitivity to near-singular orbital frames (e.g.,
circular orbits where argument of perigee is undefined) introduced
edge-case fragility.  We adopted a pragmatic compromise: a 72$\times$72
true-anomaly grid proxy (5{,}184 evaluations) with a 2.0\,km reject
threshold, deployed in shadow mode.  The sampled proxy is robust to all
orbit types and provides a reliable pre-filter ahead of the expensive
RK4 refinement stages.

\subsection{All-Pairs Fail-Open: Why Not Tiered Fallback?}

When propagation fails for an object (producing NaN or divergent states),
the engine must decide how to screen that object.  We investigated a
tiered fallback (e.g., use last-known elements for a reduced broad-phase
check) but proved that the broad phase cannot safely pair crossing-orbit
objects when elements are stale.  The only guaranteed zero-false-negative
strategy is to check the failed object against \emph{all} debris via the
full narrow-phase pipeline.  This costs more compute but eliminates the
risk of silently missing a conjunction --- consistent with the safety-first
design philosophy.

%% ═══════════════════════════════════════════════════════════════════════════
\section{Numerical Methods}\label{sec:methods}

\subsection{Orbital Propagation}

The propagator implements an adaptive three-mode scheme:

\paragraph{Mode 1 --- J2 Secular Fast Path.}
For low-eccentricity ($e \leq 0.02$), high-perigee ($\geq 500$\,km) objects
with $\Delta t \leq 30$\,s, we apply closed-form J2 secular rates to mean
elements:
%
\begin{align}
  \dot{\Omega} &= -\frac{3}{2} n J_2 \!\left(\frac{R_E}{p}\right)^{\!2} \cos i \\
  \dot{\omega} &= \frac{3}{4} n J_2 \!\left(\frac{R_E}{p}\right)^{\!2} (5\cos^{2}\!i - 1)
\end{align}
%
where $n = \sqrt{\mu / a^3}$, $p = a(1-e^2)$, and $R_E = 6378.137$\,km.
An extended fast lane admits $e \leq 0.003$, perigee $\geq 650$\,km,
$\Delta t \leq 45$\,s.  Long horizons are sub-stepped at 300\,s intervals.

\paragraph{Mode 2 --- Probe-Then-Escalate.}
Objects outside the fast lane first receive a J2-secular prediction, then a
coarse RK4 probe (max step 120\,s) validates the result.  If drift exceeds
0.5\,km in position or 0.5\,m/s in velocity, the engine escalates to full RK4.

\paragraph{Mode 3 --- Full RK4+J2.}
Forced for $\Delta t > 21{,}600$\,s (6\,h), perigee $< 350$\,km, or
$e \geq 0.98$.  Classical fourth-order Runge--Kutta with maximum substep
60\,s, integrating the acceleration:
%
\begin{equation}
  \vec{a}_{J2} = \frac{3}{2}\frac{J_2\,\mu\,R_E^2}{r^5}
  \begin{bmatrix}
    x(5z^2/r^2 - 1) \\
    y(5z^2/r^2 - 1) \\
    z(5z^2/r^2 - 3)
  \end{bmatrix}
\label{eq:j2}
\end{equation}
%
Finite-value guards after every substep prevent NaN propagation.

\paragraph{Validation.}  Maximum position error at 86{,}400\,s
(a full day) is \textbf{$<$\,0.061\,km}, verified by regression gate
\texttt{phase2\_regression\_gate}.

\paragraph{Why not RKF45 or Dormand--Prince?}
Higher-order adaptive integrators would improve accuracy per step, but
the collision threshold (100\,m) is orders of magnitude larger than the
RK4 truncation error at 60\,s substeps.  RK4 avoids the overhead of
error estimation and step rejection logic, keeping the per-object
cost predictable --- critical for deterministic tick budgets.  The
probe-then-escalate scheme already provides adaptivity at the
architectural level.

\subsection{Broad-Phase Candidate Generation}

The broad phase reduces the $O(N^2)$ all-pairs space through three
sequential filters:

\begin{enumerate}[nosep,leftmargin=*]
  \item \textbf{SMA--Inclination Band Indexing.}
    Objects are binned by semi-major axis ($\Delta a = 500$\,km) and
    inclination ($\Delta i = 20^{\circ}$).  Flat sorted arrays with binary search
    replace hash maps for cache locality.  Neighbor search extends
    $\pm 2$~bins in each dimension.

  \item \textbf{Shell Overlap Test.}
    Perigee $r_{\min} = a(1-e)$ and apogee $r_{\max} = a(1+e)$ are
    compared with a 50\,km margin.  Objects with invalid elements fall back
    to instantaneous radius $\pm 200$\,km --- an intentional
    over-approximation that preserves the zero-false-negative guarantee.

  \item \textbf{D-Criterion Gate.}
    A simplified Drummond metric:
    %
    \begin{equation}
      D = \sqrt{(\Delta a / \bar{a})^2 + \Delta e^2
          + (2\sin(\Delta i / 2))^2}
    \end{equation}
    %
    with reject threshold $D > 2.0$.  Currently deployed in
    \emph{shadow mode}: the gate logs hypothetical rejections but does
    not discard pairs, enabling safe validation before activation.
\end{enumerate}

\noindent
\textbf{Fail-Open Path.}  Objects with invalid elements,
$e > 0.2$, or radial span exceeding bin coverage bypass all cheap filters
and are checked against every debris object.  Detailed fail-open counters
are exposed in \texttt{/api/status?details=1}.

\subsection{Narrow-Phase Collision Screening}

Candidate pairs from the broad phase enter a multi-gate narrow pipeline:

\begin{enumerate}[nosep,leftmargin=*]
  \item \textbf{Plane/Phase Angular Gate.}
    Angular momentum vectors computed from $(i, \Omega)$.
    Pairs with orbital plane angle $> 75^{\circ}$ or mean anomaly phase
    separation $> 150^{\circ}$ are rejected.  Objects with $e > 0.2$ bypass
    this gate (fail-open).

  \item \textbf{MOID Proxy.}
    A sampled orbit--orbit minimum distance using a $72 \times 72$ true-anomaly
    grid (5{,}184 evaluations per pair).  Reject threshold: 2.0\,km.
    Deployed in shadow mode for safe rollout validation.

  \item \textbf{Linear TCA Approximation.}
    Four simultaneous linear-segment minimum-distance estimates
    (start velocities, end velocities backward, averaged, and secant).
    The minimum $d^2$ across all four approximations determines whether
    the pair passes.

  \item \textbf{RK4 Micro-Refinement.}
    Pairs within 0.10\,km above the screening threshold receive a
    fine-grained RK4 propagation ($\Delta t_{\max} = 5$\,s).

  \item \textbf{Full-Window RK4 Refinement.}
    Pairs within 0.20\,km receive a budget-limited, OpenMP-parallelized
    full-window refinement (16~samples, 1\,s substep).
    Thread-safe atomic budget decrement prevents compute runaway:
    base budget 64, dynamically scaled by pair count and step size,
    clamped to $[8, 192]$.
\end{enumerate}

\noindent
Uncertainty promotion adds 0.10\,km to the screening band for pairs with
relative speed exceeding 8\,km/s.  The effective screening distance is:
%
\begin{equation}
  d_{\text{screen}} = \underbrace{0.100}_{\text{collision}} +
                      \underbrace{0.020}_{\text{TCA guard}} = 0.120 \;\text{km}
\end{equation}

\subsection{Maneuver Planning}

\subsubsection{Autonomous COLA}

When a predictive CDM record reaches CRITICAL severity ($d < 100$\,m),
the engine triggers an autonomous evasion burn search over a discrete grid:

\begin{itemize}[nosep,leftmargin=*]
  \item 4~burn-time fractions of the remaining lead time
  \item 4~$\Delta V$ scales: $\{0.5, 1, 2, 4\} \times 0.001$\,km/s
  \item 10~RTN directions: $\pm T$, $\pm R$, 4~diagonal $T/R$, $\pm N$
\end{itemize}

The RTN frame is constructed as:
\begin{equation}
  \hat{r} = \frac{\vec{r}}{|\vec{r}|}, \quad
  \hat{n} = \frac{\vec{r} \times \vec{v}}{|\vec{r} \times \vec{v}|}, \quad
  \hat{t} = \hat{n} \times \hat{r}
\end{equation}

Candidates are ranked by a multi-criteria comparator: safety class
(miss distance relative to target 0.15\,km) $\to$ $|\Delta V|$ $\to$
normal component magnitude $\to$ slot error $\to$ predicted miss distance.
The top 3~critical records per satellite are evaluated.

\subsubsection{Fuel Accounting}

Propellant expenditure follows the Tsiolkovsky relation:
%
\begin{equation}
  \Delta m = m_{\text{current}} \!\left(1 - e^{-|\Delta\vec{v}| / v_e}\right),
  \quad v_e = I_{\text{sp}} \cdot g_0
\label{eq:tsiolkovsky}
\end{equation}
%
with $I_{\text{sp}} = 300$\,s, $g_0 = 9.80665$\,m/s$^2$,
$m_{\text{dry}} = 500$\,kg, $m_{\text{fuel,init}} = 50$\,kg.
A 5\%~fuel guard (2.5\,kg) triggers automatic graveyard orbit transfer.

\subsubsection{Constraint Enforcement}

\begin{itemize}[nosep,leftmargin=*]
  \item \textbf{Cooldown}: 600\,s mandatory rest between burns on the same
    satellite.
  \item \textbf{Thrust cap}: $|\Delta V| \leq 15$\,m/s per burn.
  \item \textbf{LOS validation}: Upload windows scanned in 20\,s steps
    across a 1{,}800\,s horizon against 6~ground stations. 10\,s signal
    latency enforced.
  \item \textbf{Blackout handling}: Burns must be computed and queued
    before the satellite leaves the last ground-station footprint.
\end{itemize}

\subsection{Recovery Planning (CW/ZEM)}

Every evasion maneuver triggers a paired recovery burn computed via the
Clohessy--Wiltshire (CW) state-transition formulation.  The default solver
(\texttt{CW\_ZEM\_EQUIVALENT}) inverts the in-plane $2\times2$ $\Phi_{rv}$
submatrix:
%
\begin{equation}
  \Delta\vec{v}_{\text{rec}} = \alpha\,\Phi_{rv}^{-1}\,\delta\vec{r}
    + \beta\,\delta\vec{v}
\end{equation}
%
with $\alpha = 0.70$ (position correction weight),
$\beta = 0.30$ (velocity damping weight),
CW horizon $= 0.25 \times T_{\text{orbit}}$ clamped to $[300, 5400]$\,s.
The correction is capped at 15\%~of remaining $\Delta V$ budget and 40\%~of
heuristic norm to prevent over-correction.  A heuristic fallback
(proportional RTN gains) activates when the CW matrix is near-singular
($nt \approx k \cdot 2\pi$).

The station-keeping box is defined as a 10\,km spherical radius around the
nominal orbital slot.  Slot references are bootstrapped from first telemetry
and J2-secularly propagated.  The regression gate
\texttt{recovery\_slot\_gate} enforces $\Delta_{\text{slot}} \leq 0.08$
deterministically.

%% ═══════════════════════════════════════════════════════════════════════════
\section{Frontend Mission Console}\label{sec:frontend}

The frontend is a React 18 / TypeScript single-page application built with
Vite and rendered via Canvas~2D and Three.js~WebGL.  It polls the engine's
REST API at 1--2\,Hz and presents six operational views:

\begin{enumerate}[nosep,leftmargin=*]
  \item \textbf{Command Center} (\texttt{\#/command}):
    Fleet-wide posture summary with satellite counts by status,
    active CDM warnings, recent burn history, and engine health.

  \item \textbf{Ground Track} (\texttt{\#/track}):
    Mercator-projection canvas rendering all satellite markers
    with 90-minute historical trails, 90-minute forecast paths,
    live solar terminator overlay, ground station markers, and a
    persistent map legend with epoch timestamp.

  \item \textbf{Threat Assessment} (\texttt{\#/threat}):
    Polar ``bullseye'' plot of debris in the selected satellite's
    proximity.  Radial distance encodes time-to-closest-approach;
    angular position encodes approach vector.  Color-coded by
    CDM severity (green/yellow/red).

  \item \textbf{Burn Operations} (\texttt{\#/burn-ops}):
    Maneuver timeline showing evasion and recovery burns with
    blackout zone overlays and cooldown friction markers.
    Fuel gauges per satellite.

  \item \textbf{Fleet Status} (\texttt{\#/fleet-status}):
    Aggregated constellation health with fuel heatmaps,
    $\Delta V$-cost versus collisions-avoided plots, and
    station-keeping metrics.

  \item \textbf{Scorecard} (\texttt{\#/scorecard}):
    Judge-facing evidence dashboard rolling up all six PS~\S7
    criteria with weighted readiness percentage, live evidence
    matrix, and submission gap tracking.
\end{enumerate}

\noindent
Rendering uses requestAnimationFrame with Canvas~2D for the ground track and
Three.js for the 3D polar view.  The snapshot endpoint delivers geodetic
coordinates and a compressed debris tuple array (\texttt{[id, lat, lon, alt]})
to minimize payload size and maintain 60\,FPS at 10{,}000+ debris markers.

%% ═══════════════════════════════════════════════════════════════════════════
\section{Validation and Evidence}\label{sec:validation}

\subsection{Test Infrastructure}

CASCADE maintains 8 CTest gate tests, each a standalone C++ binary in
\texttt{tools/} driven by shell harnesses in \texttt{scripts/}.  All gates
are enforced in GitHub Actions CI, and the Docker build itself fails if any
test fails.

\begin{table}[h]
\centering\small
\begin{tabularx}{\linewidth}{@{}lX@{}}
\toprule
\textbf{Gate} & \textbf{Invariant} \\
\midrule
\texttt{api\_contract\_gate}       & All 5 endpoints conform to PS schema \\
\texttt{broad\_phase\_validation}  & Shell filter $\supseteq$ brute-force baseline \\
\texttt{narrow\_phase\_false\_negative} & Zero missed collisions vs.\ RK4 reference \\
\texttt{narrow\_phase\_calibration}     & Screening thresholds within tolerance \\
\texttt{phase2\_regression}        & Propagator drift $<$ 0.061\,km at 86{,}400\,s \\
\texttt{maneuver\_ops\_invariants} & Cooldown, fuel, thrust-cap enforcement \\
\texttt{recovery\_planner\_invariants} & CW/ZEM convergence, slot error $\leq 0.08$ \\
\texttt{recovery\_slot\_gate}       & Deterministic slot recovery \\
\bottomrule
\end{tabularx}
\caption{CTest safety gates and their enforced invariants.}
\label{tab:gates}
\end{table}

\subsection{Real-Data Validation}

Beyond synthetic test suites, CASCADE was validated against a
real Space-Track OMM catalog (\texttt{data.txt}, 34\,MB).  The scenario
generator (\texttt{tools/real\_data\_scenario\_gen}) parsed 30{,}000+
objects via simdjson, filtered to LEO (periapsis $< 2{,}000$\,km altitude),
classified payloads as satellites and the remainder as debris, and ran
simulation ticks at full scale.  Result: \textbf{27{,}191} LEO objects
loaded in $<$\,1\,s with zero ECI conversion failures, and 10~collisions
detected against the real catalog (2.1\,s/tick at full 27K scale).

\subsection{Safety Score Evidence}

\begin{itemize}[nosep,leftmargin=*]
  \item \textbf{False negatives}: 0 across 6 test families (broad-phase
    sanity, narrow-phase canary, false-negative gate, calibration probe,
    regression, maneuver-ops).
  \item \textbf{Fail-open counters}: Tracked per tick. Every object with
    uncertain elements is escalated, never discarded.
  \item \textbf{Screening threshold}: Effective $0.120$\,km with
    uncertainty promotion to $0.220$\,km for hypervelocity pairs.
\end{itemize}

\subsection{Fuel Efficiency Evidence}

\begin{itemize}[nosep,leftmargin=*]
  \item Tsiolkovsky fuel accounting validated in
    \texttt{maneuver\_ops\_invariants\_gate}.
  \item CW/ZEM recovery solver minimizes $|\Delta V|$ relative to
    the heuristic proportional controller.
  \item Avoidance fuel vs.\ total fuel tracked separately in burn summaries.
\end{itemize}

\subsection{Algorithmic Speed Evidence}

\begin{table}[h]
\centering\small
\begin{tabular}{@{}lrl@{}}
\toprule
\textbf{Scenario} & \textbf{Mean Tick} & \textbf{Scale} \\
\midrule
50 sats / 10{,}000 debris   & 13.3\,ms & Baseline \\
100 sats / 20{,}000 debris  & 72.6\,ms & 2$\times$ stress \\
\bottomrule
\end{tabular}
\caption{Tick latency benchmarks (\texttt{phase3\_tick\_benchmark}).}
\label{tab:perf}
\end{table}

\noindent
The broad phase reduces candidate count from $O(N \cdot M)$ to a small
fraction via SMA/inclination binning.  OpenMP parallelization of the
narrow-phase sweep yields a 31.5\% speedup at 2$\times$ scale.
At full real-catalog scale (27{,}191 LEO objects), tick latency is
2.1\,s --- dominated by the sheer pair count but still well within
the interactive demonstration budget.

\paragraph{Performance optimizations.}
Hot-path profiling identified \texttt{std::pow(rmag, 5.0)} in the J2
acceleration as a 10--50$\times$ bottleneck versus explicit multiplication
(\texttt{r2*r2*rmag}).  Precomputing shared J2 secular terms
($\cos i$, $\sin i$, $a^{3.5}$) across the three rate functions
eliminated redundant transcendentals.  The MOID sample count was
increased from 24 to 72 after analysis showed that
$15^{\circ}$~resolution left 1{,}780\,km arc gaps at LEO altitudes,
rendering the gate effectively a no-op.

\subsection{Deployment Verification}

The Docker API smoke gate (\texttt{scripts/docker\_api\_smoke.sh}) performs
an automated release-verification cycle:

\begin{enumerate}[nosep,leftmargin=*]
  \item Build Docker image from \texttt{Dockerfile}
  \item Boot container on auto-selected port
  \item Probe \texttt{/api/status}, POST telemetry, POST simulate/step,
    GET snapshot
  \item Validate each response with Python assertions
  \item Tear down container
\end{enumerate}

\noindent
This gate runs in CI and catches regressions before any push reaches main.

%% ═══════════════════════════════════════════════════════════════════════════
\section{Deployment}\label{sec:deploy}

\subsection{Docker Build}

CASCADE uses a 3-stage multi-stage Dockerfile:

\begin{table}[h]
\centering\small
\begin{tabularx}{\linewidth}{@{}llX@{}}
\toprule
\textbf{Stage} & \textbf{Base} & \textbf{Purpose} \\
\midrule
builder  & \texttt{ubuntu:22.04} & GCC-12, CMake, CTest gates \\
frontend & \texttt{node:24-slim}  & \texttt{npm ci} + Vite production build \\
runtime  & \texttt{ubuntu:22.04} & Minimal: \texttt{libgomp1} only \\
\bottomrule
\end{tabularx}
\caption{Docker build stages.}
\label{tab:docker}
\end{table}

\noindent
BuildKit cache mounts accelerate rebuilds (\texttt{apt}, \texttt{\_deps},
\texttt{sccache}).  A \texttt{PROJECTBONK\_PREFETCH\_ONLY} step pre-fetches
all FetchContent dependencies before source copy, maximizing layer-cache hits.
The runtime image runs as a non-root user (\texttt{bonk}).

\subsection{Quick Start}

\begin{lstlisting}[language=bash]
# Build and run
docker build -t cascade:local .
docker run --rm -p 8000:8000 cascade:local

# Smoke test
curl -s http://localhost:8000/api/status \
  | python3 -m json.tool

# Full verification
./scripts/docker_api_smoke.sh
./scripts/run_ready_demo.sh
\end{lstlisting}

\subsection{CI Pipeline}

GitHub Actions runs on every push to \texttt{main}:

\begin{enumerate}[nosep,leftmargin=*]
  \item CMake configure + build (\texttt{gcc-12}, Release)
  \item CTest: all 8 safety gates
  \item Frontend: \texttt{npm ci \&\& npm run build}
  \item Docker API smoke gate
\end{enumerate}

%% ═══════════════════════════════════════════════════════════════════════════
\section{Limitations and Future Work}

\begin{itemize}[nosep,leftmargin=*]
  \item \textbf{Propagation fidelity}:
    The engine uses J2-only perturbation.  Atmospheric drag is applied
    heuristically for perigee $< 350$\,km orbits.  Lunar/solar
    third-body and SRP effects are not modeled.  For the hackathon's
    days-level prediction horizon, J2 dominates all other perturbations
    by 2--3 orders of magnitude, so this is an acceptable simplification.

  \item \textbf{MOID solver}:
    The current 72$\times$72 sampled proxy achieves sufficient accuracy
    for screening but lacks the sub-meter precision of a Gronchi~(2005)
    analytical quartic solver.  We implemented the analytical path but
    encountered edge-case fragility on near-circular orbits (undefined
    $\omega$); the sampled proxy is robust to all orbit types.

  \item \textbf{Maneuver optimization}:
    The grid-search COLA planner explores a discrete 160-candidate space
    (4~times $\times$ 4~$\Delta V$ $\times$ 10~directions).
    A continuous NLP or sequential convex programming formulation would
    yield lower-$\Delta V$ evasion sequences, but the discrete search
    completes in $<$\,1\,ms and gives deterministic, auditable results.

  \item \textbf{Dense-scene FPS}:
    Frontend rendering at 10{,}000+ debris has been tested in development
    but formal 60\,FPS proof under worst-case load needs finalization.

  \item \textbf{Filter activation}:
    The D-criterion, plane/phase, and MOID gates are deployed in shadow
    mode.  Activating them would reduce narrow-phase compute by an
    estimated 40--60\% but requires additional calibration against
    adversarial edge cases to maintain the zero-false-negative guarantee.
\end{itemize}

%% ═══════════════════════════════════════════════════════════════════════════
\section*{Acknowledgments}

We thank the organizers of the National Space Hackathon 2026 at IIT~Delhi
for designing a challenging and technically rigorous problem statement.
Space-Track (\texttt{space-track.org}) provided the real OMM catalog used
for full-scale validation.

%% ═══════════════════════════════════════════════════════════════════════════
\appendix

\section{Key Physical Constants}\label{app:constants}

\begin{table}[h]
\centering\small
\begin{tabular}{@{}llr@{}}
\toprule
\textbf{Symbol} & \textbf{Description} & \textbf{Value} \\
\midrule
$\mu$       & Earth grav.\ parameter  & 398\,600.4418\,km$^3$/s$^2$ \\
$R_E$       & Earth equatorial radius  & 6\,378.137\,km \\
$J_2$       & Oblateness coefficient   & $1.08263 \times 10^{-3}$ \\
$d_{\text{coll}}$ & Collision threshold & 0.100\,km \\
$I_{\text{sp}}$ & Specific impulse      & 300\,s \\
$g_0$       & Standard gravity         & 9.80665\,m/s$^2$ \\
$m_{\text{dry}}$ & Satellite dry mass   & 500\,kg \\
$m_{\text{fuel}}$ & Initial propellant  & 50\,kg \\
$|\Delta V|_{\max}$ & Thrust limit     & 15\,m/s \\
$\tau_{\text{cd}}$ & Cooldown period    & 600\,s \\
$m_{\text{EOL}}$ & Fuel EOL guard       & 2.5\,kg \\
$\tau_{\text{sig}}$ & Signal latency    & 10\,s \\
$r_{\text{box}}$ & Station-keeping box  & 10\,km \\
\bottomrule
\end{tabular}
\caption{System constants (from \texttt{src/types.hpp}).}
\label{tab:constants}
\end{table}

\section{CDM Severity Tiers}\label{app:cdm}

\begin{table}[h]
\centering\small
\begin{tabular}{@{}llc@{}}
\toprule
\textbf{Tier} & \textbf{Condition} & \textbf{Auto-COLA} \\
\midrule
CRITICAL & $d < 0.100$\,km  & Yes \\
WARNING  & $d < 1.0$\,km    & No (monitor) \\
WATCH    & $d < 5.0$\,km    & No (log)     \\
\bottomrule
\end{tabular}
\caption{Conjunction severity classification.}
\label{tab:cdm}
\end{table}

\section{Verification Commands}\label{app:cmds}

\begin{lstlisting}[language=bash]
# Build (local)
cmake -S . -B build -DCMAKE_BUILD_TYPE=Release
cmake --build build --target ProjectBONK

# Run all safety gates
ctest --test-dir build --output-on-failure

# Individual gates
./scripts/api_contract_gate.sh ./build
./scripts/broad_phase_sanity_gate.sh
./scripts/narrow_phase_false_negative_gate.sh
./scripts/recovery_planner_invariants_gate.sh

# Docker release verification
./scripts/docker_api_smoke.sh

# Frontend build
cd frontend && npm ci && npm run build

# Ready demo
./scripts/run_ready_demo.sh
\end{lstlisting}

\section{NSGA-II Tuner Parameters}\label{app:tuner}

The offline multi-objective tuner implements the NSGA-II algorithm
(Deb~et~al., 2002) via a Julia/\texttt{jluna} bridge.  The optimizer
drives 69~parameters across 3~tiers:

\begin{itemize}[nosep,leftmargin=*]
  \item \textbf{Tier 1 --- Broad phase} (9~params): SMA bin width,
    inclination bin width, shell margin, neighbor search range,
    D-criterion threshold, fail-open eccentricity cutoff.
  \item \textbf{Tier 2 --- Narrow phase} (21~params): plane/phase
    angular thresholds, MOID grid resolution and reject threshold,
    TCA guard margin, micro-refinement band, full-refine budget
    tiers, uncertainty promotion velocity cutoff.
  \item \textbf{Tier 3 --- Maneuver planner} (39~params): COLA
    $\Delta V$ scales, RTN direction weights, CW horizon bounds,
    position/velocity correction blend, recovery budget cap,
    CDM scanner substep and horizon.
\end{itemize}

\noindent
Three objectives are minimized simultaneously: \emph{safety risk}
(number of false negatives; any $>0$ renders the solution infeasible
via feasibility-first constraint handling), \emph{fuel consumption}
(total $\Delta V$ across the constellation), and \emph{tick latency}
(wall-clock time for one simulation step).  Verified end-to-end:
20~population $\times$ 5~generations (120~evaluations), all candidates
safe (safety\_risk~$= 0$), best tick 13.3\,ms from 23\,ms baseline
defaults.

All~69 constants are overridable at runtime via \texttt{PROJECTBONK\_*}
environment variables, enabling parameter sweeps without recompilation.

\section{Development Timeline}\label{app:timeline}

The project was developed over approximately 4~weeks across 60+~commits.
Table~\ref{tab:timeline} highlights the major milestones.

\begin{table}[h]
\centering\small
\begin{tabularx}{\linewidth}{@{}lX@{}}
\toprule
\textbf{Phase} & \textbf{Milestone} \\
\midrule
Week~1 & Telemetry ingestion, SoA state store, J2 propagator, Kepler solver, basic REST endpoints \\
Week~2 & Conservative broad-phase (shell + D-criterion), narrow-phase pipeline (TCA + MOID + RK4 refine), maneuver math (CW/ZEM), fuel and cooldown logic \\
Week~3 & Shadow-gate observability, NSGA-II tuner, real-data scenario generator (27K objects), OpenMP parallelization, 69-param env-override system \\
Week~4 & Frontend dashboard (6~pages), 3D$\to$2D pivot, Docker multi-stage build, CI gates, PS alignment, technical report \\
\bottomrule
\end{tabularx}
\caption{Development timeline and major milestones.}
\label{tab:timeline}
\end{table}

\end{document}
